\subsection{Fourier Transformation S. 23}

La trasformata di Fourier mostra \textbf{quali frequenze} ci sono in un segnale
e \textbf{quanto sono forti}.

\[
\underbrace{x(t)}_{\text{segnale nel tempo}}
\quad \circ\!\!-\!\!\bullet \quad
\underbrace{X(\omega)}_{\text{segnale in frequenza}}
\]

\begin{itemize}
  \item $x(t)$: com'è fatto il segnale \textbf{nel tempo} (come lo vedi sull'oscilloscopio).
  \item $X(\omega)$: lo \textbf{spettro}, cioè quanto c'è di ogni frequenza $\omega$ nel segnale
        (come lo vedi sull'analizzatore di spettro).
\end{itemize}

Die Fourier Transformation erweitert die spektrale Analyse auf \textbf{deterministische Energiesignale}. Herleitung siehe Skript.

Damit die Fourier Transformierte existiert muss die Funktion $x(t)$ absolut integrierbar ($\int_{-\infty}^{\infty} |x(t)| dt \leq \infty$) sein, sowie deren Fourier Transformierte $X(\omega)$ (sonst kann man sie nicht mehr zurück transformieren). Ist $f(t)$ eine Linearkombination von absolut integrierbaren Funktionen wie $x(t)$ und die resultierenden $F(\omega)$ von Funktionen wie $X(\omega)$ , dann ist $f(t)$ ebenfalls transformierbar (Linearität).
\begin{table}[H]
	\begin{tabularx}{\textwidth}{X X}
		\begin{minipage}{\linewidth}
			Fouriertransformierte
			\begin{equation*}
				X(\omega) = \int_{-\infty}^{\infty} x(t) \cdot e^{-j\omega t}dt
			\end{equation*}
		\end{minipage}
		&
		\begin{minipage}{\linewidth}
			Rücktransformierte
			\begin{equation*}
				x(t) = \frac{1}{2\pi}\int_{-\infty}^{\infty} X(\omega) \cdot e^{-j\omega t}d\omega
			\end{equation*}
	\end{minipage}
	\end{tabularx}
\end{table}

\subsubsection{Wichtige Formel}
\[e^{\pm j\omega t} = \cos(\omega t) \pm j\sin(\omega t) \qquad 
\cos(\omega t) = \frac{e^{j\omega t} + e^{-j\omega t}}{2} \qquad
\sin(\omega t) = \frac{e^{j\omega t} - e^{-j\omega t}}{2j}\]


\subsubsection{Eigenschaften der Fourier Transformation S.24}
\renewcommand{\arraystretch}{2}
\begin{tabular}{|p{2.5cm}|p{6cm}|p{8.5cm}|}
	\hline
	\multicolumn{2}{|l}{Linearität} & $\alpha\cdot f(t) + \beta\cdot g(t) \; \laplace \; \alpha\cdot F(\omega) +\beta\cdot G(\omega) \qquad (\alpha, \beta \in \mathbb{C})$  \\ 
	\hline
	\multicolumn{2}{|l}{Zeitumkehrung (Spiegelung an der Y-Achse)} & $f(-t) \; \laplace \; F(-\omega)$ \\
	\hline
	\multicolumn{2}{|l}{"Ahnlichkeit / (Zeitskalierung)}             &      		$\begin{aligned}f(\alpha t) \; &\laplace \; \frac{1}{|\alpha|}F \left (\frac{\omega}{\alpha} \right)
		\quad\alpha \in\mathbb{R}\setminus \{0\}\\
		F(\alpha\omega) \; &\Laplace \; f\left(\frac{t}{\alpha}\right)\end{aligned}$\\ 
	\hline
	\multicolumn{2}{|l}{Vertauschungssatz (Dualität)}              & 
	\begin{math}
		\begin{aligned}
			f(t) \; &\laplace \; F(\omega)\nonumber \\ 
			F(t) \; &\laplace \; 2\pi \cdot f(-\omega)
		\end{aligned}
	\end{math}\\ 
	\hline
	\hline
	\multirow{2}{*}{Verschiebung} & Zeitbereich & $f(t\pm t_0) \; \laplace \; F(\omega)e^{\pm j\omega t_0}$ \\ 
	\cline{2-3} 
	& Frequenzbereich (Modulationstheorem) & $f(t)e^{\pm j\omega_0 t} \; \laplace \; F(\omega\mp\omega_0)$ \\ 
	\hline
	\hline
	\multirow{2}{*}{Faltung} & Zeitbereich & $f(t) \ast g(t) = \int\limits_{-\infty}^{\infty} f(\tau)g(t-\tau)d\tau \; \laplace \;
	F(\omega) \cdot G(\omega)$ \\ 
	\cline{2-3} 
	& Frequenzbereich & $f(t) \cdot g(t) \; \laplace \; \frac{1}{2\pi}F(\omega) \ast G(\omega)$ \\ 
	\hline
	\hline
	\multirow{2}{*}{Ableitung} & Zeitbereich & $\frac{\partial^n f(t)}{\partial t^n} \; \laplace \; (j\omega)^n F(\omega) \qquad (n \in \mathbb{N}_0)$ \\ 
	\cline{2-3} 
	& Frequenzbereich &  $t^n f(t) \; \laplace \; j^n \frac{\partial F(\omega)}{\partial \omega^n}$\\ 
	\hline
	\hline
	\multicolumn{2}{|l}{Integration im Zeitbereich} & $\int\limits_{-\infty}^{t}f(\tau)d\tau \; \laplace \;
	\frac{F(\omega)}{j\omega}+F(0)\pi\delta(\omega)$\\ 
	\hline
	\hline
	\multicolumn{2}{|l}{Modulation} & 
	\begin{math}
		\begin{aligned}
			\cos(\alpha t) \cdot f(t)  \; &\laplace \;  \frac{1}{2}\cdot
			\left[F(\omega-\alpha) + F(\omega+\alpha)\right]\\
			\sin(\alpha t) \cdot f(t) \; &\laplace \; \frac{1}{2j}\cdot \left[F(\omega-\alpha) - F(\omega+\alpha)\right]
		\end{aligned}
	\end{math}	\\ 
	\hline
	\hline
	\multirow{2}{*}{Theorem} & Parseval's & 			 			$\int\limits_{-\infty}^{\infty}f(t)g^{\ast}(t)dt = \frac{1}{2\pi}
	\int\limits_{-\infty}^{\infty}F(\omega)G^{\ast}(\omega)d\omega$\\ 
	\cline{2-3} 
	& Bessel's & $\int\limits_{-\infty}^{\infty}|f(t)|^2 dt = \frac{1}{2\pi}
	\int\limits_{-\infty}^{\infty}|F(\omega)|^2 d\omega$ \\ 
	\hline
	\hline
	\multicolumn{2}{|l}{Anfangswerte} & 
	\begin{math}
		\begin{aligned}
			f(0)&=\frac{1}{2\pi}\int\limits_{-\infty}^{\infty}F(\omega)d\omega
			\hspace*{1cm}\\
			F(0)&=\int\limits_{-\infty}^{\infty}f(t)dt
		\end{aligned}
	\end{math} \\
	\hline
	\hline
	\multicolumn{2}{|l}{\begin{minipage}{.5\linewidth}unendlich lange Folge von $\delta$-Impulsen (im Zeit und Frequenzbereich)\end{minipage}} 
	& $\sum\limits_{n=-\infty}^{\infty} \delta(t-n\cdot t_0) \; \laplace \;
	\sum\limits_{n=-\infty}^{\infty} \frac{2\pi}{t_0}\cdot\delta\left(\omega-n\cdot
	\frac{2\pi}{t_0}\right)$ \\ 
	\hline
\end{tabular}

\subsubsection{Wichtige Fourier-Korrespondenzen}
\begin{tabular}{|l|c|c|}
\hline
\textbf{Funktion} & $\boldsymbol{x(t)}$ & $\boldsymbol{X(\omega)}$ \\
\hline
Dirac-Impuls & $\delta(t)$ & $1$ \\
\hline
Verschobener Dirac & $\delta(t-t_0)$ & $e^{-j\omega t_0}$ \\
\hline
Konstante & $1$ & $2\pi\,\delta(\omega)$ \\
\hline
Komplexe Exponentialfunktion & $e^{j\omega_0 t}$ & $2\pi\,\delta(\omega-\omega_0)$ \\
\hline
Cosinus & $\cos(\omega_0 t)$ & $\pi\left[\delta(\omega-\omega_0)+\delta(\omega+\omega_0)\right]$ \\
\hline
Sinus & $\sin(\omega_0 t)$ & $-j\pi\left[\delta(\omega-\omega_0)-\delta(\omega+\omega_0)\right]$ \\
\hline
Sprungfunktion & $u(t)$ & $\pi\,\delta(\omega)+\dfrac{1}{j\omega}$ \\
\hline
Signumfunktion & $\mathrm{sgn}(t)$ & $\dfrac{2}{j\omega}$ \\
\hline
Rechteckimpuls (Breite $2a$) & $p_a(t)$ & $\dfrac{2\sin(a\omega)}{\omega}$ \\
\hline
Rechteckimpuls (Breite $\tau$, Höhe $A$) & $A\cdot\mathrm{rect}\!\left(\dfrac{t}{\tau}\right)$ & $A\tau\cdot\dfrac{\sin(\omega\tau/2)}{\omega\tau/2}$ \\
\hline
Dreieckimpuls (Basis $2a$) & $\Lambda_a(t)$ & $\dfrac{4\sin^2(a\omega/2)}{a\,\omega^2} = a\left[\dfrac{\sin(a\omega/2)}{a\omega/2}\right]^2$ \\
\hline
Sincfunktion & $\mathrm{sinc}_a(t)=\dfrac{\sin(at)}{t}$ & $\pi\,p_a(\omega)$ \\
\hline
Sinc$^2$-Funktion & $\dfrac{\sin^2(at/2)}{t^2}$ & $\dfrac{\pi a}{2}\,\Lambda_a(\omega)$ \\
\hline
Einseitige Exponentialfunktion & $e^{-at}\,u(t),\ a>0$ & $\dfrac{1}{a+j\omega}$ \\
\hline
 & $t\,e^{-at}\,u(t),\ a>0$ & $\dfrac{1}{(a+j\omega)^2}$ \\
\hline
Zweiseitige Exponentialfunktion & $e^{-a|t|},\ a>0$ & $\dfrac{2a}{a^2+\omega^2}$ \\
\hline
Gedämpfter Cosinus & $e^{-at}\cos(\omega_0 t)\,u(t)$ & $\dfrac{a+j\omega}{(a+j\omega)^2+\omega_0^2}$ \\
\hline
Gedämpfter Sinus & $e^{-at}\sin(\omega_0 t)\,u(t)$ & $\dfrac{\omega_0}{(a+j\omega)^2+\omega_0^2}$ \\
\hline
Gaussimpuls & $e^{-at^2},\ a>0$ & $\sqrt{\dfrac{\pi}{a}}\;e^{-\omega^2/(4a)}$ \\
\hline
\end{tabular}
\renewcommand{\arraystretch}{\arraystretchOriginal}
\medbreak

\begin{minipage}{8cm}
	Periodische Signale\vspace{-1em}
	\boxedeq{%
	\sum_{n=-\infty}^{\infty}	c_n \cdot e^{jn\omega_0t} \;\; \Laplace \;\; 2\pi \sum_{n=-\infty}^{\infty} c_n \delta(\omega - n\omega_0)
	}
\end{minipage}
\medbreak

\begin{minipage}{15cm}
	Spektrum Einton Winkelmodulation\vspace{-1em}
	\boxedeq{%
	\sum_{n=-\infty}^{\infty} c_n \cdot cos((\omega_c + n\omega_m) \cdot t) \;\; \Laplace \;\;
	\sum_{n=-\infty}^{\infty}c_n\pi(\delta(\omega - (\omega_c+n\omega_m))+\delta(\omega+(\omega_c+n\omega_m)))
	}
\end{minipage}
\medbreak

Für weitere Korrespondenzpaare siehe Skript Seite 28 und \textbf{Seite 253 ff}.
